% --- scholar LaTeX fragment --------------------------------------------
% Generated by: scholar sessions export -f table
% Session: screenlit-q2
% \input this file from the body of a document whose
% preamble loads:
%   \usepackage{longtable}
%   \usepackage{booktabs}
%   \usepackage{array}
% ----------------------------------------------------------------------

\begingroup\footnotesize
\begin{longtable}{@{}%
  >{\raggedright\arraybackslash}p{0.44\textwidth}c%
  >{\raggedright\arraybackslash}p{0.13\textwidth}%
  >{\raggedright\arraybackslash}p{0.26\textwidth}@{}}
\caption{Records bearing on the claim, Question 2, abbreviated searches and screens (16 cited, 7 supporting and 17 qualifying or contradicting it, out of 3,922 unique hits; the 82 adjacent hits, the 670 false hits and 3,130 unclassified are not listed). The last column gives the reason for each inclusion; for rows decided by an agent the agent's confidence is given. A row is cited when the text cites the work, whoever decided on its records. Databases searched: OpenAlex (OA; not open access), IEEE Xplore (IEEE), Web of Science (WoS), Scopus, Semantic Scholar (S2), arXiv. Searched 2026-11-23.}\label{tab:screenlit-q2}\\
\toprule
Title & Year & Database & Reason \\
\midrule
\endfirsthead
\multicolumn{4}{@{}l}{\emph{\tablename~\thetable{} (cont.)}}\\
\toprule
Title & Year & Database & Reason \\
\midrule
\endhead
\midrule \multicolumn{4}{r@{}}{\emph{continued on next page}}\\
\endfoot
\bottomrule
\endlastfoot
Abbreviated and comprehensive literature searches led to identical or very similar effect estimates: a meta-epidemiological study & 2020 & WoS, Scopus, OA, S2 & \emph{cited}: supports the claim (0.80) \\
An evaluation of DistillerSR's machine learning-based prioritization tool for title/abstract screening -- impact on reviewer-relevant outcomes & 2020 & OA, WoS, Scopus, S2 & \emph{cited}, pending, unclassified \\
Assessing the coverage of PubMed, Embase, OpenAlex, and Semantic Scholar for automated single-database searches in living guideline evidence surveillance: a case study of the international polycystic ovary syndrome guidelines 2023 & 2025 & Scopus & \emph{cited}, pending, unclassified \\
Assessing the validity of a rapid review against a systematic literature review. A comparison of systematic literature reviews done by Cochrane with rapid reviews and the impact on meta-analyses results & 2024 & OA, WoS, Scopus, S2 & \emph{cited}: qualifies or contradicts the claim (0.70) \\
Combining abbreviated literature searches with single-reviewer screening: three case studies of rapid reviews & 2020 & OA, WoS, Scopus, S2 & \emph{cited}: qualifies or contradicts the claim (0.75) \\
Comparing the Efficacy of Rapid Review With a Systematic Review in the Software Engineering Field & 2024 & OA, WoS, Scopus, S2 & \emph{cited}: qualifies or contradicts the claim (0.70) \\
Does a rapid review version of a large epidemiological systematic review fail to identify many eligible studies, and what implications does this have for the results of the review? & 2019 & OA, S2 & \emph{cited}: qualifies or contradicts the claim (0.60) \\
Evaluating the proportion of retrieved citations and relevant studies when database selection is restricted to MEDLINE, CENTRAL, and Embase: part 4 of a meta-research study of rapid review methodological trade-offs & 2026 & OA, WoS, Scopus, S2 & \emph{cited}: qualifies or contradicts the claim (0.70) \\
Implications of applying methodological shortcuts to expedite systematic reviews: three case studies using systematic reviews from agrifood public health & 2016 & OA, S2, WoS, Scopus & \emph{cited}: qualifies or contradicts the claim (0.75) \\
Meta-epidemiologic analysis indicates that MEDLINE searches are sufficient for diagnostic test accuracy systematic reviews & 2014 & OA, WoS, Scopus, S2 & \emph{cited}: supports the claim (0.70) \\
Rapid reviews may produce different results to systematic reviews: a meta-epidemiological study & 2018 & OA, S2, WoS, Scopus & \emph{cited}: qualifies or contradicts the claim (0.80) \\
Single-reviewer abstract screening missed 13 percent of relevant studies: a crowd-based, randomized controlled trial & 2020 & OA, WoS, Scopus, S2 & \emph{cited}: qualifies or contradicts the claim (0.75) \\
The contribution of databases to the results of systematic reviews: a cross-sectional study & 2016 & OA, WoS, Scopus, S2 & \emph{cited}: supports the claim (0.70) \\
The Impact Of Searching Fewer Databases In Health Technology Assessment Rapid Reviews & 2017 & OA, S2 & \emph{cited}: supports the claim (0.65) \\
The value of a second reviewer for study selection in systematic reviews & 2019 & OA, S2, WoS, Scopus & \emph{cited}: qualifies or contradicts the claim (0.70) \\
Using data sources beyond PubMed has a modest impact on the results of systematic reviews of therapeutic interventions & 2015 & Scopus & \emph{cited}: supports the claim (0.75) \\
Empirical Testing of Alternative Search Methods to Retrieve Utility Values for Health Economic Modelling & 2024 & OA, WoS, Scopus, S2 & supports the claim (0.50) \\
Excluding non-English publications from evidence-syntheses did not change conclusions: a meta-epidemiological study & 2019 & OA, S2, WoS, Scopus & supports the claim (0.60) \\
Increasing the efficiency of study selection for systematic reviews using prioritization tools and a single-screening approach & 2023 & OA, Scopus, WoS, S2 & supports the claim (0.55) \\
Limited search strategies were effective in finding relevant nonrandomized studies & 2006 & Scopus, S2, OA, WoS & supports the claim (0.50) \\
Limits in the search date for rapid reviews of diagnostic test accuracy studies & 2022 & OA, WoS, Scopus, S2 & supports the claim (0.60) \\
Simulated changes in systematic review screening sensitivity and specificity associated with varying co-screening proportions & 2026 & OA, Scopus, S2 & supports the claim (0.55) \\
Successful incorporation of single reviewer assessments during systematic review screening: development and validation of sensitivity and work-saved of an algorithm that considers exclusion criteria and count & 2021 & OA, WoS, Scopus, S2 & supports the claim (0.55) \\
A retrospective comparison of systematic reviews with same-topic rapid reviews & 2017 & OA, WoS, Scopus, S2 & qualifies or contradicts the claim (0.60) \\
Characteristics and recovery methods of studies falsely excluded during literature screening---a systematic review & 2022 & OA, WoS, Scopus, S2 & qualifies or contradicts the claim (0.45) \\
Cochrane systematic reviews in acupuncture: Methodological diversity in database searching & 2005 & WoS, OA, S2, Scopus & qualifies or contradicts the claim (0.45) \\
Comparison of a full systematic review versus rapid review approaches to assess a newborn screening test for tyrosinemia type 1 & 2017 & OA, WoS, Scopus, S2 & qualifies or contradicts the claim (0.60) \\
Database selection in systematic reviews: an insight through clinical neurology & 2017 & S2, OA, WoS, Scopus & qualifies or contradicts the claim (0.50) \\
Error rates of human reviewers during abstract screening in systematic reviews & 2020 & OA, WoS, Scopus, S2 & qualifies or contradicts the claim (0.55) \\
Evaluating the accuracy and impact of methodological trade-offs used in rapid reviews on results and certainty of the evidence: meta-research studies in systematic reviews of chronic low back pain & 2026 & OA & qualifies or contradicts the claim (0.65) \\
Evaluating the accuracy of full text and serial screening by one vs two independent reviewers: part 2 of a meta-research study of rapid review methodological trade-offs & 2026 & OA, WoS, Scopus, S2 & qualifies or contradicts the claim (0.60) \\
Evaluating the accuracy of screening titles and abstracts for progressing to full-text screening by one vs two independent reviewers: part 1 of a meta-research study of rapid review methodological trade-offs & 2026 & OA, WoS, Scopus, S2 & qualifies or contradicts the claim (0.65) \\
Evaluating the impact of methodological trade-offs on meta-analyses and certainty of the evidence according to GRADE: part 5 of a meta-research study of rapid reviews & 2026 & OA, WoS, Scopus, S2 & qualifies or contradicts the claim (0.60) \\
Examining screening test accuracy in Tyrosinemia type 1: A comparison of rapid and systematic review Methods & 2016 & OA, S2 & qualifies or contradicts the claim (0.50) \\
Identification of randomized controlled trials in systematic reviews: accuracy and reliability of screening records & 2002 & OA, WoS, Scopus, S2 & qualifies or contradicts the claim (0.60) \\
Lessons for search strategies from a systematic review, in The Cochrane Library, of nutritional supplementation trials in patients after hip fracture & 2001 & WoS, OA, Scopus, S2 & qualifies or contradicts the claim (0.55) \\
Literature search in systematic reviews: How much is good enough? & 2023 & OA, WoS, Scopus & qualifies or contradicts the claim (0.60) \\
Literature searching for social science systematic reviews: consideration of a range of search techniques & 2010 & OA, WoS, Scopus, S2 & qualifies or contradicts the claim (0.50) \\
Network meta-analyses could be improved by searching more sources and by involving a librarian & 2014 & OA, S2, WoS, Scopus & qualifies or contradicts the claim (0.45) \\
Sources of evidence for systematic reviews of interventions in diabetes & 2005 & WoS, OA, Scopus & qualifies or contradicts the claim (0.55) \\
\end{longtable}
\endgroup
